\chapter{Eosinophilic Esophagitis}
\index{Eosinophilic Esophagitis}
\label{sec:Eosinophilic Esophagitis}

\section{Overview}
Eosinophilic esophagitis (EoE) is a chronic, immune-mediated inflammatory condition of the esophagus characterized by symptoms of esophageal dysfunction and histologic evidence of eosinophil-predominant inflammation. It affects approximately 1 in 2,000 individuals, with a strong male predominance and peak incidence in children and adults aged 20--40.
\section{Classification}

\begin{description}[font=\normalfont\bfseries,leftmargin=3em,labelwidth=2.5em]
  \item[Cause] Allergic/Immune
  \item[Organ System] Gastrointestinal/Esophagus
  \item[Pathophysiology] Th2-mediated allergic response to food or environmental allergens -> eosinophilic infiltration and degranulation in esophageal mucosa -> epithelial injury, fibrosis, and stricture formation
  \item[Duration] Chronic
  \item[Onset] Gradual
  \item[Mode of Transmission] Non-communicable (immune-mediated)
  \item[Clinical Course] Chronic, relapsing-remitting
  \item[Severity] Mild to Severe
  \item[Extent of Disease] Localized (esophagus)
  \item[Age Group] Children and Adults
  \item[Epidemiology] Male predominance (3:1); increasing incidence in developed countries
  \item[ICD-11 Code] DC20.0
\end{description}

\section{Causes}
Eosinophilic esophagitis arises from a type 2 helper T-cell (Th2)-driven allergic inflammatory response to food allergens (most commonly milk, wheat, egg, soy, and peanuts) or aeroallergens in genetically susceptible individuals. The resulting eosinophil infiltration into the esophageal epithelium releases cytotoxic granule proteins that damage tissue, promote fibrosis, and cause dysmotility. This condition develops from a combination of genetic and environmental factors that disrupt normal body processes. When these normal processes are disrupted, it leads to the symptoms and signs of the condition. Several factors can work together to trigger or make the condition worse. Knowing the specific cause helps your doctor choose the most effective treatment.

\section{Risk Factors}

\begin{itemize}[noitemsep]
  \item Atopic conditions (asthma, allergic rhinitis, atopic dermatitis, food allergies)
  \item Male sex (3:1 male-to-female ratio)
  \item Family history of eosinophilic esophagitis
  \item Living in a cold or dry climate (environmental aeroallergen exposure)
  \item Season of birth (spring/summer linked to higher risk)
\end{itemize}

\section{Signs and Symptoms}

\subsection{Early symptoms}
\begin{itemize}[noitemsep]
  \item Early manifestations of eosinophilic esophagitis may be subtle and nonspecific
\end{itemize}

\subsection{Common symptoms}
\begin{itemize}[noitemsep]
  \item \subsection{Early symptoms}
  \item \subsection{Common symptoms}
  \item \textbf{Common symptoms:}
  \item In adults: dysphagia (difficulty swallowing) and food impaction, often with solid foods
  \item In children: abdominal pain, vomiting, regurgitation, nausea, and failure to thrive
  \item Heartburn or chest pain that does not respond to proton pump inhibitors
  \item The esophagus may appear with rings (trachealization), longitudinal furrows, white exudates, or strictures on endoscopy
  \item Pain or discomfort in the affected area
  \item Difficulty performing daily activities
\end{itemize}

\subsection{Less common symptoms}
\begin{itemize}[noitemsep]
  \item Atypical or less common presentations of eosinophilic esophagitis may occur
\end{itemize}

\subsection{Emergency warning signs}
\begin{itemize}[noitemsep]
  \item Severe or worsening symptoms of eosinophilic esophagitis require immediate medical evaluation
\end{itemize}
\section{Progression and Stages}
Untreated eosinophilic esophagitis progresses from an inflammatory phenotype (edema, exudates, furrows) to a fibrostenotic phenotype (rings, strictures, narrowing) over years, driven by chronic eosinophilic inflammation. Disease activity fluctuates with allergen exposure, and progression is faster in patients with persistent untreated inflammation.

\section{Complications}

\begin{itemize}[noitemsep]
  \item Esophageal food impaction requiring emergency endoscopic removal
  \item Esophageal stricture formation requiring dilation
  \item Spontaneous esophageal perforation (Boerhaave syndrome, rare)
  \item Reduced quality of life
  \item Emotional and psychological effects
\end{itemize}

\section{Diagnosis}

\begin{itemize}[noitemsep]
  \item Upper endoscopy with esophageal biopsies showing >=15 eosinophils per high-power field after an 8-week PPI trial to exclude PPI-responsive esophageal eosinophilia
  \item Clinical history of esophageal dysfunction (dysphagia, food impaction)
  \item Lab tests to check how your organs are working
  \item Imaging tests
\end{itemize}

\section{Prevention}

\begin{itemize}[noitemsep]
  \item Avoidance of identified trigger foods (guided by allergy testing or elimination diets)
  \item Management of concurrent atopic conditions
  \item Go for regular check-ups so any problems are caught early
  \item Follow your doctor's screening recommendations
\end{itemize}

\section{Treatment Options}

\begin{itemize}[noitemsep]
  \item Topical swallowed corticosteroids (fluticasone or budesonide slurry) to reduce esophageal eosinophilic inflammation
  \item Dietary elimination therapy (six-food elimination diet, targeted elimination based on allergy testing, or elemental diet)
  \item Proton pump inhibitors as first-line therapy for patients with PPI-responsive esophageal eosinophilia
  \item Lifestyle changes and self-care
  \item Surgery when needed (endoscopic dilation for strictures)
\end{itemize}

\section{Living with It}

\begin{itemize}[noitemsep]
  \item Follow your treatment plan as prescribed
  \item Keep up with regular appointments with your healthcare provider
  \item Adjust your daily habits as needed
  \item Reach out to family, friends, or support groups for help
  \item Keep learning about your condition and how to manage it
\end{itemize}

\section{Outlook}
Eosinophilic esophagitis is a chronic condition requiring long-term management, but most patients achieve symptom control with dietary therapy or topical corticosteroids. Your outlook depends on the specific condition, how advanced it is when found, and how well it responds to treatment. Getting treated early generally leads to better results. Many people manage this condition effectively with the right treatment. Regular check-ups are important to monitor your condition and adjust treatment as needed.
